\documentclass[12pt]{article} % Клас документа, шрифт 12pt

\usepackage[utf8]{inputenc} % UTF-8 кодування
\usepackage[T2A]{fontenc} % Шрифти кирилиці
\usepackage[ukrainian]{babel} % Українська мова
\usepackage{geometry} % Поля сторінки
\geometry{a4paper, margin=2cm} % Формат A4, поля 2 см

\usepackage{xcolor} % Кольори
\usepackage{tcolorbox} % Кольорові рамки
\usepackage{graphicx} % Зображення
\usepackage{amsmath, amsthm, amssymb} % Математика, теореми, символи

\usepackage{caption} % Підписи до рисунків
\captionsetup{labelsep=period} % Роздільник у підписах "Рис. 1."

\newtheorem{theorem}{Теорема} % Середовище для теорем

\usepackage{fancyhdr} % Колонтитули
\pagestyle{fancy} % Використання fancyhdr
\fancyhf{} % Очищення стандартних колонтитулів
\setlength{\headheight}{15pt} % Висота верхнього колонтитулу

\newcommand{\student}{Левенець Владислава} % ПІБ студента
\newcommand{\labdate}{\today} % Дата

\fancyhead[L]{\student} % Лівий верхній колонтитул
\fancyhead[R]{\labdate} % Правий верхній колонтитул
\fancyfoot[C]{\thepage} % Нижній колонтитул — номер сторінки

\begin{document}

\begin{center}
    \Large \textbf{Лабораторна робота № 9} \\
    \large \textbf{Тема: Набір математичних текстів та автоматизація в LaTeX}
\end{center}

\section*{Мета роботи}
\textit{Навчитися вводити математичні формули та структурувати документ, використовувати кольорові рамки для визначень, теорем та прикладів, а також вставляти зображення}.

\section{Короткі теоретичні відомості}

\begin{tcolorbox}[colback=blue!10!white, colframe=red!80!black, title=Визначення]
% Введіть визначення квадрата, змініть колір рамки
Квадрат — це чотирикутник з усіма сторонами рівними та всіма кутами прямими.
\end{tcolorbox}

% Введіть текст \\На рисунку ... дивись зразок. 
На рисунку 1 зображено квадрат ABCD. Прямокутник є паралелограмом, тому і квадрат - паралелограм, у якого всі сторони рівні, тобто він є і ромбом. Отже, квадрат має в властиості прямокутника та ромба.

%додати зображення, файл foto1.png 
\begin{figure}[h]
    \centering
    \includegraphics[scale=1]{foto1.png}
    \caption{Діагоналі квадрата}
    \label{fig:foto1}
\end{figure}

\medskip
\textit{Властивості квадрата:}
\begin{enumerate}
    \item Усі кути квадрата прямі.
    \item Периметр $P_{ABCD} = 4AB$.
    \item Діагоналі квадрата між собою рівні (\ref{fig:foto2})
    \item Діагоналі квадрата взаємно перпендинелярні й точкою перетину діляться навпіл (\ref{fig:foto2})
    \item Діагоналі квадрата ділять його куи навпіл, тобто утворюють кути $45^\circ$ зі сторонами квадрата(\ref{fig:foto2})
    \item Точка перетину діагоналей квадрата рівновіддалена від усіх його вершин:   $AO=BO=CO=DO$
   %Введіть властивості 2-6 за зразком
   
\end{enumerate}

%додати зображення, файл foto2.png 
\begin{figure}[h]
    \centering
    \includegraphics[scale=1]{foto2.png}
    \caption{Діагоналі квадрата}
    \label{fig:foto2}
\end{figure}

\begin{tcolorbox}[colback=green!5!white, colframe=green!75!black, title=Приклад 1]

Точки $K$ і $M$ належать відповідно діагоналям $BD$ і $AC$ квадрата $ABCD$, причому 
$AM = \dfrac{1}{4}AC,\quad BK = \dfrac{1}{4}BD$. 
Доведіть, що $\triangle ADM = \triangle BAK$~(\ref{fig:foto3}).

\end{tcolorbox}

%додати зображення, файл foto3.png 
\begin{figure}[h]
    \centering
    \includegraphics[scale=1]{foto3.png}
    \caption{Діагоналі квадрата}
    \label{fig:foto3}
\end{figure}

\begin{proof}
\begin{enumerate}
    \item $\angle MAD = \angle ABK = 45^\circ$, а $AD = AB$ як сторони квадрата.
    
    \item Оскільки $AC = BD$ як діагоналі квадрата і 
    $AM = \dfrac{1}{4}AC,\quad BK = \dfrac{1}{4}BD$, то $AM = BK$.
    
    \item Отже, $\triangle AMD = \triangle BKA$ 
    (за двома сторонами та кутом між ними, перша ознака рівності трикутників).
    
    \item Звідси випливає рівність відповідних елементів трикутників.
\end{enumerate}
\end{proof}

\section{Ознаки квадрата}

\begin{tcolorbox}[colback=blue!5!white, colframe=blue!75!black, title=Ознаки квадрата]
\begin{theorem}
1) Якщо діагоналі прямокутника взаємно перпендикулярні, то він є квадратом \\
2) Якщо діагоналі ромба між собою рівні то він є квадратом
\end{theorem}
\end{tcolorbox}

\begin{proof}
%Введіть текст доведення використовуючи зразок
1) Прямокутник є паралелограмом, а паралелограм зі взаємно перпендикулярними діагоналями є ромбом. Отже, у заданого прямокутника всі сторони рівні, а тому він є квадратом.
 \\
2) Ромб є паралелограмом, а паралелограм з рівними діагоналями є прямокутником. Отже, у розглянутому ромбі всі кути прямі, а тому він є квадратом
\end{proof}


\end{document}